\chapter{Dilaton--radion action и dimensional transmutation}

Absolute-scale no-go показывает, что классический parent action должен либо
содержать физическую шкалу, либо порождать её квантово. Рассмотрим сначала
минимальный classically scale-invariant сектор.

\section{Однополевой классический no-go}

Для дилатона \(\chi\) единственный ренормируемый потенциал без размерных
коэффициентов имеет вид
\begin{equation}
 V_{\mathrm{cl}}(\chi)=\lambda_\chi\chi^4.
\end{equation}
При \(\lambda_\chi>0\) единственный минимум равен \(\chi=0\); при
\(\lambda_\chi=0\) всё направление плоское; при
\(\lambda_\chi<0\) потенциал не ограничен снизу. Ненулевой устойчивый
масштаб классически не возникает.

\section{Два scale-поля}

Пусть \(\rho\) --- radion-модуль той же массовой размерности. Потенциал
\begin{equation}
 V_{\mathrm{ratio}}
 =\lambda\left(\rho^2-\alpha\chi^2\right)^2
\end{equation}
фиксирует отношение
\begin{equation}
 \frac{\rho}{\chi}=\sqrt\alpha,
\end{equation}
но сохраняет общее преобразование
\begin{equation}
 (\chi,\rho)\longmapsto t(\chi,\rho).
\end{equation}
В точке \((\chi,\rho)=(1,\sqrt\alpha)\) Hessian имеет собственные значения
\begin{equation}
 8\alpha\lambda(\alpha+1),
 \qquad 0.
\end{equation}
Нулевая мода является неизбежным классическим дилатоном.

\begin{theorem}[Classical dilaton flat-direction no-go]
Полиномиальный classically scale-invariant dilaton--radion potential может
фиксировать безразмерные отношения vacuum expectation values, но не
абсолютный общий масштаб. Любой ненулевой классический vacuum образует
непрерывную scale-орбиту.
\end{theorem}

\section{Radius-inversion кандидат}

Функционал
\begin{equation}
 V_{\mathrm{dual}}(\varphi)
 =A\left(e^{2\varphi}+e^{-2\varphi}-2\right),
 \qquad
 R=R_0e^\varphi,
\end{equation}
имеет единственный минимум \(R=R_0\). Но равенство коэффициентов двух
экспонент требует симметрии
\begin{equation}
 \varphi\longmapsto-\varphi,
 \qquad
 R\longmapsto\frac{R_0^2}{R}.
\end{equation}
Обычный momentum-only спектр поля на окружности такой radius inversion не
реализует. Для неё необходимы winding states, dual field или отдельно
доказанная modular symmetry.

\begin{nogocriterion}[Unproved radius duality]
Потенциал \(V_{\mathrm{dual}}\) запрещено использовать для фиксации
\(R=R_0\), пока radius inversion не выведена из объявленного field menu.
Иначе равенство коэффициентов является скрытой scale-настройкой.
\end{nogocriterion}

\section{Coleman--Weinberg кандидат}

Квантовая scale anomaly допускает эффективный потенциал
\begin{equation}
 V_{\mathrm{CW}}(\chi)
 =B\chi^4
 \left(
 \log\frac{\chi^2}{\mu^2}-\frac12
 \right),
 \qquad B>0.
\end{equation}
Он удовлетворяет
\begin{equation}
 V_{\mathrm{CW}}'(\mu)=0,
 \qquad
 V_{\mathrm{CW}}''(\mu)=8B\mu^2>0.
\end{equation}
Минимальная совместная конструкция имеет вид
\begin{equation}
 V_{\mathrm{joint}}(\chi,D_F)
 =
 V_{\mathrm{CW}}(\chi)
 +\chi^4\Tr_{\mathrm{orb}}
 \left(
 \frac{D_F^2}{\chi^2}-\mathbf1
 \right)^2.
\end{equation}
Её минимум даёт
\begin{equation}
 \chi=\mu,
 \qquad
 D_F^2=\mu^2\mathbf1,
 \qquad
 |x|=|z|=\frac{\mu}{\sqrt2}.
\end{equation}
При идентификации \(\chi=R^{-1}\) одновременно стабилизируется
\begin{equation}
 R_*=\mu^{-1}.
\end{equation}

\begin{proposition}[Quantum joint-vacuum candidate]
Coleman--Weinberg anomaly может одним vacuum equation связать dilaton,
radion и конечный Dirac scale. При \(B>0\) scale-direction получает
положительную массу, а finite-Dirac направления сохраняют положительный
flattening Hessian.
\end{proposition}

\section{Что ещё не выведено}

Параметр \(\mu\) в фиксированном порядке является renormalization point.
Физическая transmutation scale должна быть записана как RG-инвариант
\begin{equation}
 \Lambda_{\mathrm{DT}}
 =
 \mu
 \exp\left(
 -\int^{g(\mu)}
 \frac{dg}{\beta(g)}
 \right).
\end{equation}
Её численное значение определяется beta-функциями и граничным значением
хотя бы одной безразмерной связи. Поэтому dimensional transmutation
заменяет явный mass parameter на quantum boundary condition; она не
предсказывает число без полного RG-модуля.

В текущей версии не вычислены:
\begin{enumerate}
 \item коэффициент \(B\) из полного KO6 finite module и KK-башен;
 \item знак \(B\) после ghost, gauge и fermion contributions;
 \item RG boundary condition для dimensionless couplings;
 \item backreaction дилатона на factor measures и orbit half-trace.
\end{enumerate}

\begin{nogocriterion}[Quantum scale gate]
Dilaton--radion branch считается закрытой положительно только если один
пререгистрированный spectrum вычисляет \(B>0\), beta-функции и
\(\Lambda_{\mathrm{DT}}\) без использования лептонных масс как boundary
condition. В противном случае \(\mu\) остаётся переименованным scale input.
\end{nogocriterion}

\section{Следующий вычислительный шаг}

Следующий аудит должен построить spectrum ledger полного кандидата и
вычислить supertrace коэффициент Coleman--Weinberg:
\begin{equation}
 B\propto
 \Str M^4(\chi)/\chi^4.
\end{equation}
Первый kill-критерий прост: если \(B\le0\), минимальный quantum vacuum
неустойчив и dilaton--radion маршрут закрывается в текущем field menu.

Машинный аудит сохранён в
\path{s2t_v3_dilaton_radion_transmutation_results.json}.